\documentclass[a4paper,12pt]{article}
\usepackage[utf8]{inputenc}
\usepackage[T1]{fontenc}
\usepackage[french]{babel}
\usepackage{geometry}
\usepackage{array}
\usepackage{booktabs}
\usepackage{xcolor}
\usepackage{colortbl}
\usepackage{longtable}
\usepackage{hyperref}
\usepackage{parskip}
\usepackage{graphicx}
\usepackage{fancyhdr}
\usepackage{titlesec}

\geometry{margin=2cm, top=2.5cm, bottom=2.5cm}

% Couleurs CapAvenir
\definecolor{capblue}{RGB}{0, 87, 184}      % #0057B8
\definecolor{capdark}{RGB}{0, 59, 142}      % #003B8E
\definecolor{caporange}{RGB}{255, 106, 0}   % #FF6A00
\definecolor{rowgray}{RGB}{242, 246, 252}
\definecolor{rowwhite}{RGB}{255, 255, 255}

% En-tête et pied de page
\pagestyle{fancy}
\fancyhf{}
\fancyhead[L]{\small\color{capdark}\textbf{CapAvenir} --- Plateforme d'Orientation Intelligente}
\fancyhead[R]{\small\color{capdark}Cahier des Cas d'Utilisation}
\fancyfoot[C]{\small\thepage}
\renewcommand{\headrulewidth}{1pt}
\renewcommand{\headrule}{\hbox to\headwidth{\color{capblue}\leaders\hrule height \headrulewidth\hfill}}

% Style des titres de section
\titleformat{\section}{\Large\bfseries\color{capdark}}{}{0em}{}[\titlerule]
\titleformat{\subsection}{\large\bfseries\color{capblue}}{}{0em}{}

% Commande en-tête de cas d'utilisation
\newcommand{\ucheader}[2]{%
  \noindent\colorbox{capblue}{\parbox{\dimexpr\linewidth-2\fboxsep}{%
    \color{white}\bfseries\large Cas d'utilisation~: #1 \hfill \normalsize Acteur~: #2
  }}\\[2pt]
}

% Colonnes de tableau
\newcolumntype{L}{>{\raggedright\arraybackslash}p{3.8cm}}
\newcolumntype{R}{>{\raggedright\arraybackslash}p{11.2cm}}

\hypersetup{
    colorlinks=true,
    linkcolor=capblue,
    urlcolor=caporange,
    pdfauthor={Projet PFE --- CapAvenir},
    pdftitle={Cahier des Cas d'Utilisation --- CapAvenir}
}

%=============================================================
\begin{document}
%=============================================================

\begin{titlepage}
  \centering
  \vspace*{2cm}
  {\color{capblue}\rule{\linewidth}{3pt}}\\[0.5em]
  {\Huge\bfseries\color{capdark} Cahier des Cas d'Utilisation}\\[0.6em]
  {\large\color{capblue} Plateforme CapAvenir --- Système SIAEPI d'Orientation Intelligente}\\[0.3em]
  {\color{caporange}\rule{\linewidth}{1.5pt}}\\[1.5em]
  {\large Spécification fonctionnelle complète}\\[0.3em]
  {\normalsize 4 acteurs $\cdot$ 28 cas d'utilisation documentés}\\[2em]
  \begin{tabular}{ll}
    \textbf{Projet}     & Fin d'Études (PFE) \\
    \textbf{Plateforme} & CapAvenir v5.0 \\
    \textbf{Stack}      & Laravel 10 $\cdot$ PHP 8.2 $\cdot$ Moteur SIAEPI $\cdot$ CAT 2PL \\
    \textbf{Année}      & 2025 -- 2026 \\
  \end{tabular}
  \vfill
  {\color{capblue}\rule{\linewidth}{3pt}}
\end{titlepage}

\tableofcontents
\newpage

%=============================================================
\section{Présentation de la plateforme CapAvenir}
%=============================================================

\textbf{CapAvenir} est une plateforme web d'orientation académique intelligente développée avec le framework \textbf{Laravel 10} (PHP 8.2). Elle intègre le moteur psychométrique \textbf{SIAEPI v5.0}, un moteur de test adaptatif (\textit{CAT --- Computerized Adaptive Testing}) basé sur le modèle IRT 2PL, ainsi que plusieurs modules d'aide à la décision destinés aux étudiants bacheliers tunisiens.

\subsection*{Architecture et acteurs}
La plateforme distingue quatre acteurs principaux~:
\begin{enumerate}
  \item \textbf{Visiteur / Invité} : utilisateur non authentifié accédant à la landing page publique.
  \item \textbf{Étudiant} : bénéficiaire principal du système d'orientation.
  \item \textbf{Conseiller d'Orientation} : encadreur supervisant les dossiers psychométriques des étudiants.
  \item \textbf{Administrateur} : gestionnaire de la configuration globale de la plateforme.
\end{enumerate}

\subsection*{Modules fonctionnels implémentés}
\begin{itemize}
  \item \textbf{Test CAT adaptatif} (RIASEC 6 dimensions + GATB 8 aptitudes + Big Five) via \texttt{AdaptiveTestEngine} et \texttt{TestManager}
  \item \textbf{Moteur de recommandation SIAEPI} avec score pondéré, filtres d'exclusion durs et optimisation de Pareto
  \item \textbf{Simulateur What-If} (scénarios prospectifs via \texttt{FutureSimulatorService})
  \item \textbf{Comparateur de filières} multi-critères (jusqu'à 3 filières en parallèle)
  \item \textbf{CV Builder} avec export PDF / DOCX
  \item \textbf{Assistant Nova} (chatbot IA via API Gemini avec stratégie de fallback)
  \item \textbf{Pipeline d'orientation} 3 étapes (test $\rightarrow$ score $\rightarrow$ recommandations)
  \item \textbf{Tableau de bord CRM Conseiller} avec analyse prédictive et historique de suivi
  \item \textbf{Panneau Admin} (CRUD utilisateurs, questions RIASEC, import Excel filières, audit)
  \item \textbf{Double authentification (2FA)} par OTP email
\end{itemize}

\newpage

%=============================================================
\section{Acteur : Visiteur / Invité}
%=============================================================

Le visiteur désigne tout utilisateur non authentifié arrivant sur la plateforme CapAvenir. Il peut accéder à la landing page publique, passer un test RIASEC en mode découverte et soumettre une demande de contact.

%-------------------------------------------------------------
\subsection{S'inscrire (Créer un compte)}
%-------------------------------------------------------------
\ucheader{S'inscrire (Créer un compte)}{Visiteur}

\noindent
\begin{tabular}{|L|R|}
\hline
\rowcolor{capblue}\textcolor{white}{\textbf{Champ}} & \textcolor{white}{\textbf{Description}} \\
\hline
\rowcolor{rowgray}Identifiant & UC-V01 \\
\hline
Cas d'utilisation & S'inscrire (Créer un compte) \\
\hline
\rowcolor{rowgray}Acteur & Visiteur \\
\hline
Précondition & Le visiteur n'est pas connecté et dispose d'une adresse email valide. \\
\hline
\rowcolor{rowgray}Postcondition & Un compte utilisateur est créé en base de données avec le rôle \texttt{student} par défaut~; un profil vide est initialisé. \\
\hline
Scénario nominal &
\begin{enumerate}
  \item Le visiteur clique sur « S'inscrire » depuis la landing page.
  \item Le système affiche le formulaire d'inscription (nom, email, mot de passe, confirmation, rôle).
  \item Le visiteur saisit ses informations et soumet.
  \item Le système vérifie l'unicité de l'email via la règle \texttt{unique:users} et la robustesse du mot de passe (\texttt{min:8}).
  \item Le système enregistre l'utilisateur avec le rôle \texttt{student} et crée un profil associé.
  \item Le système redirige vers la page de connexion avec un message de succès.
\end{enumerate} \\
\hline
\rowcolor{rowgray}Scénario alternatif &
\begin{itemize}
  \item Email déjà existant $\rightarrow$ Validation Laravel retourne une erreur et invite l'utilisateur à se connecter.
  \item Données invalides (mot de passe trop court, format email incorrect) $\rightarrow$ Les champs en erreur sont surlignés en rouge.
\end{itemize} \\
\hline
Composants impliqués & \texttt{Auth/RegisterController}, \texttt{User} (Model), \texttt{Profile} (Model), vue \texttt{auth/register.blade.php} \\
\hline
\end{tabular}

\vspace{1em}

%-------------------------------------------------------------
\subsection{S'authentifier avec double facteur (2FA)}
%-------------------------------------------------------------
\ucheader{S'authentifier (Connexion + OTP 2FA)}{Visiteur}

\noindent
\begin{tabular}{|L|R|}
\hline
\rowcolor{capblue}\textcolor{white}{\textbf{Champ}} & \textcolor{white}{\textbf{Description}} \\
\hline
\rowcolor{rowgray}Identifiant & UC-V02 \\
\hline
Cas d'utilisation & S'authentifier avec double facteur (2FA) \\
\hline
\rowcolor{rowgray}Acteur & Visiteur (possédant un compte) \\
\hline
Précondition & L'utilisateur possède un compte actif et non bloqué (\texttt{is\_blocked = 0}). \\
\hline
\rowcolor{rowgray}Postcondition & Session sécurisée active~; redirection vers le tableau de bord correspondant au rôle. \\
\hline
Scénario nominal &
\begin{enumerate}
  \item L'utilisateur saisit son email et son mot de passe sur \texttt{/login}.
  \item \texttt{LoginController} vérifie les identifiants.
  \item Si le 2FA est activé, le système génère un code OTP à 6 chiffres, le stocke dans la session et l'envoie par email via \texttt{TwoFactorController}.
  \item L'utilisateur est redirigé vers \texttt{/two-factor} pour saisir le code.
  \item Le système valide le code OTP. La session est authentifiée.
  \item Le middleware \texttt{two-factor} lève le verrou et redirige selon le rôle (\texttt{student}, \texttt{counselor}, \texttt{admin}).
\end{enumerate} \\
\hline
\rowcolor{rowgray}Scénario alternatif &
\begin{itemize}
  \item Identifiants invalides $\rightarrow$ Message d'erreur, compte non ouvert.
  \item Compte bloqué (\texttt{is\_blocked = 1}) $\rightarrow$ Le système refuse la connexion avec un message explicatif.
  \item Code OTP expiré ou incorrect $\rightarrow$ Blocage de l'accès~; possibilité de renvoi via \texttt{POST /two-factor/resend}.
\end{itemize} \\
\hline
Composants impliqués & \texttt{Auth/LoginController}, \texttt{Auth/TwoFactorController}, middleware \texttt{two-factor} \\
\hline
\end{tabular}

\vspace{1em}

%-------------------------------------------------------------
\subsection{Envoyer un message de contact}
%-------------------------------------------------------------
\ucheader{Envoyer un message de contact}{Visiteur}

\noindent
\begin{tabular}{|L|R|}
\hline
\rowcolor{capblue}\textcolor{white}{\textbf{Champ}} & \textcolor{white}{\textbf{Description}} \\
\hline
\rowcolor{rowgray}Identifiant & UC-V03 \\
\hline
Cas d'utilisation & Envoyer un message de contact \\
\hline
\rowcolor{rowgray}Acteur & Visiteur \\
\hline
Précondition & Accès à la landing page publique. \\
\hline
\rowcolor{rowgray}Postcondition & Le message est enregistré dans la table \texttt{contacts} et l'administrateur peut le consulter depuis son panneau. \\
\hline
Scénario nominal &
\begin{enumerate}
  \item Le visiteur accède au formulaire de contact en bas de la landing page.
  \item Il saisit son nom, email, objet et message, puis clique sur « Envoyer ».
  \item \texttt{ContactController::store()} valide les données et enregistre le contact.
  \item Un message de succès est affiché.
\end{enumerate} \\
\hline
\rowcolor{rowgray}Scénario alternatif &
\begin{itemize}
  \item Champs vides ou email invalide $\rightarrow$ Le système bloque l'envoi et surligne les champs en erreur.
\end{itemize} \\
\hline
Composants impliqués & \texttt{ContactController}, \texttt{Contact} (Model), route publique \texttt{POST /contact} \\
\hline
\end{tabular}

\vspace{1em}

%-------------------------------------------------------------
\subsection{Passer le test RIASEC en mode invité}
%-------------------------------------------------------------
\ucheader{Passer le test RIASEC (mode découverte invité)}{Visiteur}

\noindent
\begin{tabular}{|L|R|}
\hline
\rowcolor{capblue}\textcolor{white}{\textbf{Champ}} & \textcolor{white}{\textbf{Description}} \\
\hline
\rowcolor{rowgray}Identifiant & UC-V04 \\
\hline
Cas d'utilisation & Passer le test RIASEC en mode invité \\
\hline
\rowcolor{rowgray}Acteur & Visiteur \\
\hline
Précondition & Accès à la plateforme (aucune authentification requise). \\
\hline
\rowcolor{rowgray}Postcondition & Profil RIASEC calculé et résultats affichés. Les données sont stockées en base si l'invité s'inscrit, ou purgées à expiration de session. \\
\hline
Scénario nominal &
\begin{enumerate}
  \item Le visiteur accède au test RIASEC via \texttt{GET /riasec/question}.
  \item \texttt{RiasecTestController::start()} initialise une session de test via \texttt{POST /riasec/question}.
  \item Le moteur \texttt{TestManager} génère dynamiquement la première question en fonction du modèle adaptatif 2PL.
  \item Le visiteur répond question par question~; \texttt{storeAnswer()} enregistre chaque réponse et met à jour $\theta$ (niveau latent).
  \item Le \texttt{BehavioralAnalyzer} surveille les temps de réponse et la cohérence comportementale.
  \item L'arrêt précoce (\textit{Early Stopping}) est déclenché dès que l'erreur de mesure SEM descend sous le seuil.
  \item Le profil RIASEC est calculé et affiché à l'adresse \texttt{GET /riasec/resultats}.
\end{enumerate} \\
\hline
\rowcolor{rowgray}Scénario alternatif &
\begin{itemize}
  \item Test abandonné $\rightarrow$ La session expire et les réponses temporaires sont purgées automatiquement.
\end{itemize} \\
\hline
Composants impliqués & \texttt{RiasecTestController}, \texttt{TestManager}, \texttt{AdaptiveTestEngine}, \texttt{BehavioralAnalyzer}, \texttt{EarlyStoppingService}, middleware \texttt{riasec.test} \\
\hline
\end{tabular}

\newpage

%=============================================================
\section{Acteur : Étudiant}
%=============================================================

L'étudiant est le bénéficiaire central de CapAvenir. Il utilise les outils d'intelligence artificielle et les modules d'analyse pour construire et simuler son projet d'orientation académique.

%-------------------------------------------------------------
\subsection{Passer le test d'orientation adaptatif (CAT)}
%-------------------------------------------------------------
\ucheader{Passer le test d'orientation intelligent (CAT + GATB)}{Étudiant}

\noindent
\begin{tabular}{|L|R|}
\hline
\rowcolor{capblue}\textcolor{white}{\textbf{Champ}} & \textcolor{white}{\textbf{Description}} \\
\hline
\rowcolor{rowgray}Identifiant & UC-E01 \\
\hline
Cas d'utilisation & Passer le test d'orientation adaptatif (CAT) \\
\hline
\rowcolor{rowgray}Acteur & Étudiant \\
\hline
Précondition & Étudiant authentifié, aucun test en cours d'analyse. \\
\hline
\rowcolor{rowgray}Postcondition & Profil psychométrique complet généré~: 6 dimensions RIASEC, 8 aptitudes GATB, indicateurs Big Five, score de cohérence et flags comportementaux. Profil enregistré dans \texttt{ProfileRiasec}. \\
\hline
Scénario nominal &
\begin{enumerate}
  \item L'étudiant accède au module de test depuis son tableau de bord.
  \item \texttt{TestManager} génère dynamiquement la question suivante selon l'algorithme 2PL (paramètres de difficulté $b$ et discrimination $a$).
  \item L'étudiant répond sur échelle de Likert 1--5.
  \item Le \texttt{BehavioralAnalyzer} surveille en temps réel~: temps de réponse, patterns d'incohérence, indices de fraude.
  \item \texttt{AdaptiveTestEngine} met à jour le niveau latent $\theta$ par méthode de Newton-Raphson.
  \item Les étapes 2 à 5 se répètent jusqu'au déclenchement de l'\textit{Early Stopping} (SEM $<$ seuil configuré).
  \item \texttt{GatbCalculator} calcule les 8 aptitudes cognitives (raisonnement verbal, numérique, spatial, etc.).
  \item \texttt{PostTestValidator} valide l'intégrité du profil et le stocke définitivement.
\end{enumerate} \\
\hline
\rowcolor{rowgray}Scénario alternatif &
\begin{itemize}
  \item Fraude détectée (temps de réponse anormal, cohérence $<$ seuil) $\rightarrow$ \texttt{BehavioralAnalyzer} lève un drapeau \texttt{fraud\_flag} dans \texttt{ProfileRiasec} et pénalise les scores de certitude.
  \item Test interrompu $\rightarrow$ L'état est sauvegardé en session pour reprise ultérieure.
\end{itemize} \\
\hline
Composants impliqués & \texttt{TestManager}, \texttt{AdaptiveTestEngine}, \texttt{BehavioralAnalyzer}, \texttt{EarlyStoppingService}, \texttt{GatbCalculator}, \texttt{IrtCalibrator}, \texttt{PostTestValidator}, \texttt{ProfileRiasec} (Model) \\
\hline
\end{tabular}

\vspace{1em}

%-------------------------------------------------------------
\subsection{Suivre le pipeline d'orientation (3 étapes)}
%-------------------------------------------------------------
\ucheader{Suivre le pipeline d'orientation (Profil $\rightarrow$ Test $\rightarrow$ Recommandations)}{Étudiant}

\noindent
\begin{tabular}{|L|R|}
\hline
\rowcolor{capblue}\textcolor{white}{\textbf{Champ}} & \textcolor{white}{\textbf{Description}} \\
\hline
\rowcolor{rowgray}Identifiant & UC-E02 \\
\hline
Cas d'utilisation & Suivre le pipeline d'orientation guidé \\
\hline
\rowcolor{rowgray}Acteur & Étudiant \\
\hline
Précondition & Étudiant authentifié avec rôle \texttt{student}. \\
\hline
\rowcolor{rowgray}Postcondition & Les 3 étapes du pipeline sont complétées~; l'étudiant accède à ses recommandations finales. \\
\hline
Scénario nominal &
\begin{enumerate}
  \item L'étudiant clique sur « Commencer mon orientation » depuis son tableau de bord.
  \item \textbf{Étape 1 --- Profil académique}~: l'étudiant renseigne sa section de BAC et ses notes~; \texttt{OrientationPipelineController::storeStep1()} calcule et enregistre le score FG.
  \item \textbf{Étape 2 --- Test CAT}~: redirection vers \texttt{/riasec/question} pour passer le test psychométrique adaptatif complet.
  \item \textbf{Étape 3 --- Recommandations}~: le moteur SIAEPI génère le top-12 des filières correspondantes.
\end{enumerate} \\
\hline
\rowcolor{rowgray}Scénario alternatif &
\begin{itemize}
  \item Profil incomplet lors de l'étape 1 $\rightarrow$ Le système bloque la progression et invite à compléter les champs manquants.
\end{itemize} \\
\hline
Composants impliqués & \texttt{OrientationPipelineController}, routes \texttt{student.pipeline}, vue \texttt{student/pipeline/step1\_score.blade.php} \\
\hline
\end{tabular}

\vspace{1em}

%-------------------------------------------------------------
\subsection{Consulter les recommandations personnalisées (SIAEPI)}
%-------------------------------------------------------------
\ucheader{Consulter les recommandations de filières (SIAEPI v5.0)}{Étudiant}

\noindent
\begin{tabular}{|L|R|}
\hline
\rowcolor{capblue}\textcolor{white}{\textbf{Champ}} & \textcolor{white}{\textbf{Description}} \\
\hline
\rowcolor{rowgray}Identifiant & UC-E03 \\
\hline
Cas d'utilisation & Consulter les recommandations de filières (SIAEPI v5.0) \\
\hline
\rowcolor{rowgray}Acteur & Étudiant \\
\hline
Précondition & Profil psychométrique CAT complété et profil académique renseigné (score FG calculé). \\
\hline
\rowcolor{rowgray}Postcondition & Top-12 des filières affiché avec scores de matching, Gap Analysis et indicateurs Pareto. \\
\hline
Scénario nominal &
\begin{enumerate}
  \item L'étudiant accède à l'espace « Recommandations » via \texttt{GET /recommendations}.
  \item \texttt{SiaepiRecommendationEngine} extrait les dimensions RIASEC, les aptitudes GATB, le Big Five et le score FG officiel.
  \item Le moteur applique les filtres d'exclusion durs (incompatibilité académique, seuils SDO).
  \item Le score d'adéquation pondéré global est calculé par filière (axes~: Vocation, Académique, Marché, Accessibilité).
  \item Le système identifie les filières optimales au sens de Pareto.
  \item Les recommandations, la Gap Analysis et les indices de cohérence sont affichés.
\end{enumerate} \\
\hline
\rowcolor{rowgray}Scénario alternatif &
\begin{itemize}
  \item Profil académique manquant $\rightarrow$ Le système redirige l'étudiant vers la page de saisie de ses notes BAC.
  \item Aucune filière ne passe les filtres durs $\rightarrow$ Affichage d'un message explicatif invitant l'étudiant à consulter son conseiller.
\end{itemize} \\
\hline
Composants impliqués & \texttt{SiaepiRecommendationEngine}, \texttt{ScoreFGService}, \texttt{StudentController::showRecommendations()}, \texttt{Filiere} (Model) \\
\hline
\end{tabular}

\vspace{1em}

%-------------------------------------------------------------
\subsection{Gérer son profil académique (Notes du BAC)}
%-------------------------------------------------------------
\ucheader{Gérer son profil académique (Notes et Score FG)}{Étudiant}

\noindent
\begin{tabular}{|L|R|}
\hline
\rowcolor{capblue}\textcolor{white}{\textbf{Champ}} & \textcolor{white}{\textbf{Description}} \\
\hline
\rowcolor{rowgray}Identifiant & UC-E04 \\
\hline
Cas d'utilisation & Gérer son profil académique \\
\hline
\rowcolor{rowgray}Acteur & Étudiant \\
\hline
Précondition & Étudiant authentifié. \\
\hline
\rowcolor{rowgray}Postcondition & Section de BAC et notes enregistrées~; score FG calculé automatiquement et persisté dans \texttt{Profile}. \\
\hline
Scénario nominal &
\begin{enumerate}
  \item L'étudiant ouvre la page « Mon Profil » via \texttt{GET /student/profil}.
  \item Il sélectionne sa section officielle de BAC tunisien dans la liste déroulante.
  \item Le système adapte dynamiquement les champs aux matières correspondantes.
  \item L'étudiant saisit sa moyenne générale et ses notes par matière, puis valide via \texttt{PUT /student/profil}.
  \item \texttt{StudentProfileController::update()} vérifie la validité des notes (0 à 20) et calcule le score global via la formule officielle.
  \item Le profil académique est sauvegardé.
\end{enumerate} \\
\hline
\rowcolor{rowgray}Scénario alternatif &
\begin{itemize}
  \item Note en dehors des bornes [0, 20] $\rightarrow$ Le système refuse la sauvegarde et affiche une erreur de validation.
\end{itemize} \\
\hline
Composants impliqués & \texttt{StudentProfileController}, \texttt{ScoreFGService}, \texttt{Profile} (Model), vue \texttt{student/profil.blade.php} \\
\hline
\end{tabular}

\vspace{1em}

%-------------------------------------------------------------
\subsection{Utiliser le comparateur de filières}
%-------------------------------------------------------------
\ucheader{Utiliser le comparateur de filières (multi-critères)}{Étudiant}

\noindent
\begin{tabular}{|L|R|}
\hline
\rowcolor{capblue}\textcolor{white}{\textbf{Champ}} & \textcolor{white}{\textbf{Description}} \\
\hline
\rowcolor{rowgray}Identifiant & UC-E05 \\
\hline
Cas d'utilisation & Utiliser le comparateur de filières \\
\hline
\rowcolor{rowgray}Acteur & Étudiant \\
\hline
Précondition & Étudiant authentifié. \\
\hline
\rowcolor{rowgray}Postcondition & Tableau comparatif affiché avec graphique radar interactif et indicateurs colorés. \\
\hline
Scénario nominal &
\begin{enumerate}
  \item L'étudiant accède à \texttt{GET /student/comparateur}.
  \item Il sélectionne jusqu'à 3 filières universitaires depuis la liste.
  \item Il valide via \texttt{POST /student/comparateur/data}.
  \item \texttt{ComparateurController::comparer()} extrait les données de chaque filière~: seuil d'accès SDO, taux d'employabilité, salaire moyen, adéquation RIASEC, localisation.
  \item Le système présente un tableau comparatif avec indicateurs visuels colorés et un graphique radar.
\end{enumerate} \\
\hline
\rowcolor{rowgray}Scénario alternatif &
\begin{itemize}
  \item Moins de 2 filières sélectionnées $\rightarrow$ La comparaison est bloquée avec un message d'aide.
\end{itemize} \\
\hline
Composants impliqués & \texttt{ComparateurController}, \texttt{Filiere} (Model), vue \texttt{student/comparateur/index.blade.php} \\
\hline
\end{tabular}

\vspace{1em}

%-------------------------------------------------------------
\subsection{Utiliser le simulateur de futurs (What-If)}
%-------------------------------------------------------------
\ucheader{Utiliser le simulateur de futurs (What-If)}{Étudiant}

\noindent
\begin{tabular}{|L|R|}
\hline
\rowcolor{capblue}\textcolor{white}{\textbf{Champ}} & \textcolor{white}{\textbf{Description}} \\
\hline
\rowcolor{rowgray}Identifiant & UC-E06 \\
\hline
Cas d'utilisation & Utiliser le simulateur de futurs (What-If) \\
\hline
\rowcolor{rowgray}Acteur & Étudiant \\
\hline
Précondition & Étudiant authentifié. \\
\hline
\rowcolor{rowgray}Postcondition & Scénario prospectif calculé (delta impact, nouvelles filières débloquées, ROI prévisionnel) et enregistré dans \texttt{simulation\_history}. \\
\hline
Scénario nominal &
\begin{enumerate}
  \item L'étudiant ouvre le « Simulateur de Futurs » via \texttt{GET /student/whatif}.
  \item Il sélectionne un axe de simulation~: variation de notes, coût des études, ROI à 10 ans.
  \item Il ajuste les réglettes ou saisit des valeurs hypothétiques.
  \item \texttt{WhatIfController::calculer()} ou \texttt{simulerAvance()} invoque \texttt{FutureSimulatorService} pour calculer le delta d'impact.
  \item Le système affiche~: nouveau score FG estimé, filières débloquées, salaire prévisionnel, probabilité de succès.
  \item L'étudiant sauvegarde le scénario dans son historique.
\end{enumerate} \\
\hline
\rowcolor{rowgray}Scénario alternatif &
\begin{itemize}
  \item Suppression d'un scénario historique $\rightarrow$ \texttt{DELETE /student/whatif/historique/\{simulation\}} supprime définitivement la simulation.
\end{itemize} \\
\hline
Composants impliqués & \texttt{WhatIfController}, \texttt{FutureSimulatorService}, \texttt{SimulationHistory} (Model), vues \texttt{student/whatif/} \\
\hline
\end{tabular}

\vspace{1em}

%-------------------------------------------------------------
\subsection{Interagir avec l'assistant Nova (Chatbot IA)}
%-------------------------------------------------------------
\ucheader{Interagir avec l'assistant Nova (Chatbot Gemini)}{Étudiant}

\noindent
\begin{tabular}{|L|R|}
\hline
\rowcolor{capblue}\textcolor{white}{\textbf{Champ}} & \textcolor{white}{\textbf{Description}} \\
\hline
\rowcolor{rowgray}Identifiant & UC-E07 \\
\hline
Cas d'utilisation & Interagir avec l'assistant Nova \\
\hline
\rowcolor{rowgray}Acteur & Étudiant \\
\hline
Précondition & Étudiant authentifié. \\
\hline
\rowcolor{rowgray}Postcondition & Réponse contextuelle personnalisée renvoyée via l'API Gemini (ou Ollama en fallback local). \\
\hline
Scénario nominal &
\begin{enumerate}
  \item L'étudiant accède au module Nova via \texttt{GET /student/orientation/nova}.
  \item Il formule sa question en langage naturel dans l'interface de chat.
  \item \texttt{NovaOrientationController::analyze()} enrichit la requête avec les métadonnées de l'étudiant (profil RIASEC, notes BAC, filière cible).
  \item Le système contacte l'API Gemini en appliquant la stratégie de \textit{fallback} (rotation entre modèles Gemini Flash et Pro).
  \item La réponse structurée est affichée à l'étudiant.
  \item Le résultat final est accessible via \texttt{GET /student/orientation/nova/resultat}.
\end{enumerate} \\
\hline
\rowcolor{rowgray}Scénario alternatif &
\begin{itemize}
  \item Limite de requêtes API dépassée (HTTP 429) $\rightarrow$ Le système bascule automatiquement sur le modèle suivant de la chaîne de secours (\texttt{OllamaService} en local).
\end{itemize} \\
\hline
Composants impliqués & \texttt{NovaOrientationController}, \texttt{NovaOrientationService}, \texttt{OllamaService}, vues \texttt{student/nova/} \\
\hline
\end{tabular}

\vspace{1em}

%-------------------------------------------------------------
\subsection{Créer et gérer ses CV (CV Builder)}
%-------------------------------------------------------------
\ucheader{Créer et gérer ses CV (CV Builder + Export)}{Étudiant}

\noindent
\begin{tabular}{|L|R|}
\hline
\rowcolor{capblue}\textcolor{white}{\textbf{Champ}} & \textcolor{white}{\textbf{Description}} \\
\hline
\rowcolor{rowgray}Identifiant & UC-E08 \\
\hline
Cas d'utilisation & Créer et gérer ses CV \\
\hline
\rowcolor{rowgray}Acteur & Étudiant \\
\hline
Précondition & Étudiant authentifié. \\
\hline
\rowcolor{rowgray}Postcondition & CV enregistré en base (tables \texttt{cv\_profiles}, \texttt{cv\_education}, \texttt{cv\_experiences}, \texttt{cv\_skills}, \texttt{cv\_languages})~; export PDF et DOCX disponibles. \\
\hline
Scénario nominal &
\begin{enumerate}
  \item L'étudiant se rend dans « CV Builder » via \texttt{GET /student/cv}.
  \item Il crée un nouveau CV via \texttt{GET /student/cv/create} ou édite un brouillon existant.
  \item Il renseigne ses rubriques~: coordonnées, formations, expériences, compétences, langues.
  \item Il sélectionne un template graphique (\textit{modern}).
  \item Le système affiche une prévisualisation dynamique en temps réel.
  \item Il télécharge son CV au format PDF (\texttt{PdfGeneratorService}) ou Word (\texttt{DocxGeneratorService}).
\end{enumerate} \\
\hline
\rowcolor{rowgray}Scénario alternatif &
\begin{itemize}
  \item Duplication de CV $\rightarrow$ \texttt{POST /student/cv/\{cvProfile\}/duplicate} clone le profil CV complet.
  \item Suppression $\rightarrow$ \texttt{DELETE /student/cv/\{cvProfile\}} supprime le CV et toutes ses données liées.
\end{itemize} \\
\hline
Composants impliqués & \texttt{CvBuilderController}, \texttt{PdfGeneratorService}, \texttt{DocxGeneratorService}, \texttt{CvProfile}, \texttt{CvEducation}, \texttt{CvExperience}, \texttt{CvSkill}, \texttt{CvLanguage} (Models), template \texttt{cv\_templates/modern\_pdf.blade.php} \\
\hline
\end{tabular}

\vspace{1em}

%-------------------------------------------------------------
\subsection{Gérer son portfolio de compétences}
%-------------------------------------------------------------
\ucheader{Gérer son portfolio de compétences}{Étudiant}

\noindent
\begin{tabular}{|L|R|}
\hline
\rowcolor{capblue}\textcolor{white}{\textbf{Champ}} & \textcolor{white}{\textbf{Description}} \\
\hline
\rowcolor{rowgray}Identifiant & UC-E09 \\
\hline
Cas d'utilisation & Gérer son portfolio de compétences \\
\hline
\rowcolor{rowgray}Acteur & Étudiant \\
\hline
Précondition & Étudiant authentifié. \\
\hline
\rowcolor{rowgray}Postcondition & Réalisation ou projet académique consigné dans \texttt{portfolio\_items} et rattaché au profil de l'étudiant. \\
\hline
Scénario nominal &
\begin{enumerate}
  \item L'étudiant accède à son portfolio depuis son tableau de bord.
  \item Il clique sur « Ajouter un projet ».
  \item Il saisit le titre, la description, les tags de compétences, et téléverse un justificatif.
  \item \texttt{PortfolioController::store()} valide et enregistre l'item.
  \item Le projet est rattaché au profil de l'étudiant.
\end{enumerate} \\
\hline
\rowcolor{rowgray}Scénario alternatif &
\begin{itemize}
  \item Suppression d'un item $\rightarrow$ \texttt{DELETE /student/portfolio/\{portfolio\}} supprime l'item.
\end{itemize} \\
\hline
Composants impliqués & \texttt{PortfolioController}, \texttt{PortfolioItem} (Model) \\
\hline
\end{tabular}

\vspace{1em}

%-------------------------------------------------------------
\subsection{Générer sa roadmap de carrière}
%-------------------------------------------------------------
\ucheader{Générer sa roadmap de carrière (IA Gemini)}{Étudiant}

\noindent
\begin{tabular}{|L|R|}
\hline
\rowcolor{capblue}\textcolor{white}{\textbf{Champ}} & \textcolor{white}{\textbf{Description}} \\
\hline
\rowcolor{rowgray}Identifiant & UC-E10 \\
\hline
Cas d'utilisation & Générer sa roadmap de carrière \\
\hline
\rowcolor{rowgray}Acteur & Étudiant \\
\hline
Précondition & Profil RIASEC/CAT actif et filière recommandée disponible. \\
\hline
\rowcolor{rowgray}Postcondition & Itinéraire professionnel structuré (formations, compétences, certifications, jalons) enregistré dans \texttt{career\_roadmaps}. \\
\hline
Scénario nominal &
\begin{enumerate}
  \item L'étudiant clique sur « Générer ma Roadmap » (\texttt{POST /student/roadmap}).
  \item \texttt{RoadmapController::generate()} récupère le trigramme Holland et le profil RIASEC de l'étudiant.
  \item Le service IA (Gemini via \texttt{ChatbotController}) construit un itinéraire professionnel structuré~: formations, compétences à acquérir, certifications, secteurs clés.
  \item La roadmap est enregistrée et affichée sous forme de frise interactive dans le tableau de bord.
\end{enumerate} \\
\hline
\rowcolor{rowgray}Scénario alternatif &
\begin{itemize}
  \item Service d'IA indisponible $\rightarrow$ Le système affiche un parcours générique basé sur la filière recommandée numéro 1.
\end{itemize} \\
\hline
Composants impliqués & \texttt{RoadmapController}, \texttt{ChatbotController}, \texttt{CareerRoadmap} (Model) \\
\hline
\end{tabular}

\vspace{1em}

%-------------------------------------------------------------
\subsection{Communiquer via la messagerie interne}
%-------------------------------------------------------------
\ucheader{Communiquer via la messagerie interne}{Étudiant}

\noindent
\begin{tabular}{|L|R|}
\hline
\rowcolor{capblue}\textcolor{white}{\textbf{Champ}} & \textcolor{white}{\textbf{Description}} \\
\hline
\rowcolor{rowgray}Identifiant & UC-E11 \\
\hline
Cas d'utilisation & Communiquer via la messagerie interne \\
\hline
\rowcolor{rowgray}Acteur & Étudiant \\
\hline
Précondition & Étudiant authentifié. \\
\hline
\rowcolor{rowgray}Postcondition & Message stocké dans la table \texttt{messages}~; indicateur de non-lecture mis à jour. \\
\hline
Scénario nominal &
\begin{enumerate}
  \item L'étudiant ouvre la messagerie interne depuis son tableau de bord.
  \item Il clique sur « Nouveau message » vers son conseiller assigné.
  \item Il rédige et envoie via \texttt{UserMessageController}.
  \item Le système stocke le message et met à jour le compteur de notifications.
\end{enumerate} \\
\hline
\rowcolor{rowgray}Scénario alternatif &
\begin{itemize}
  \item Message vide $\rightarrow$ Le système bloque l'envoi avec une erreur de validation.
\end{itemize} \\
\hline
Composants impliqués & \texttt{UserMessageController}, \texttt{Message} (Model) \\
\hline
\end{tabular}

\newpage

%=============================================================
\section{Acteur : Conseiller d'Orientation}
%=============================================================

Le conseiller supervise les dossiers psychométriques des étudiants via un tableau de bord CRM avancé. Il peut analyser les recommandations générées par l'IA, planifier des rendez-vous, rédiger des observations de suivi, valider les matchings et envoyer des communications omnicanales.

%-------------------------------------------------------------
\subsection{Consulter le tableau de bord analytique}
%-------------------------------------------------------------
\ucheader{Consulter le tableau de bord analytique (CRM Conseiller)}{Conseiller d'orientation}

\noindent
\begin{tabular}{|L|R|}
\hline
\rowcolor{capblue}\textcolor{white}{\textbf{Champ}} & \textcolor{white}{\textbf{Description}} \\
\hline
\rowcolor{rowgray}Identifiant & UC-C01 \\
\hline
Cas d'utilisation & Consulter le tableau de bord analytique \\
\hline
\rowcolor{rowgray}Acteur & Conseiller d'orientation \\
\hline
Précondition & Conseiller authentifié avec le rôle \texttt{counselor}. \\
\hline
\rowcolor{rowgray}Postcondition & KPIs de la cohorte, tendances d'orientation, heatmaps sectorielles et insights IA affichés. \\
\hline
Scénario nominal &
\begin{enumerate}
  \item Le conseiller accède à son tableau de bord via \texttt{GET /counselor}.
  \item \texttt{CounselorController::index()} agrège~: liste des étudiants, \texttt{cohortStats} (répartition sectorielle), \texttt{kpis} (taux de réussite, risque d'abandon, progression), \texttt{heatmaps} (filières saturées / émergentes), \texttt{iaInsights} (alertes IA explicables).
  \item Le système affiche également les données de benchmark national et international (établissements, régions, tendances mondiales).
  \item Les rendez-vous à venir du conseiller sont listés avec horodatage.
\end{enumerate} \\
\hline
\rowcolor{rowgray}Scénario alternatif &
\begin{itemize}
  \item Aucun étudiant dans la base $\rightarrow$ Affichage d'un état vide avec message d'invitation.
\end{itemize} \\
\hline
Composants impliqués & \texttt{CounselorController}, \texttt{User}, \texttt{Profile}, \texttt{Appointment} (Models), vue \texttt{counselor/dashboard.blade.php} \\
\hline
\end{tabular}

\vspace{1em}

%-------------------------------------------------------------
\subsection{Consulter le dossier complet d'un étudiant}
%-------------------------------------------------------------
\ucheader{Consulter le dossier psychométrique d'un étudiant}{Conseiller d'orientation}

\noindent
\begin{tabular}{|L|R|}
\hline
\rowcolor{capblue}\textcolor{white}{\textbf{Champ}} & \textcolor{white}{\textbf{Description}} \\
\hline
\rowcolor{rowgray}Identifiant & UC-C02 \\
\hline
Cas d'utilisation & Consulter le dossier complet d'un étudiant \\
\hline
\rowcolor{rowgray}Acteur & Conseiller d'orientation \\
\hline
Précondition & Conseiller authentifié. \\
\hline
\rowcolor{rowgray}Postcondition & Fiche complète affichée~: archétype de profil, risk flags, timeline CRM, plan de coaching, indicateurs prédictifs, historique de décisions. \\
\hline
Scénario nominal &
\begin{enumerate}
  \item Le conseiller clique sur un étudiant depuis son tableau de bord.
  \item \texttt{CounselorController::showStudent()} charge le profil complet, les tentatives de test, les rendez-vous, et génère les données CRM~:
  \begin{itemize}
    \item \texttt{crmPriority}~: niveau de priorité de suivi (Urgent, Surveillance, Standard, Haute performance).
    \item \texttt{riskFlags}~: incohérence d'orientation, motivation faible, stress élevé, faible compatibilité.
    \item \texttt{crmTimeline}~: chronologie des événements (tests, RDV, notes, simulations, rapports).
    \item \texttt{archetype} + \texttt{coachingPlan}~: archétype IA et plan de coaching automatisé.
    \item \texttt{crmPredictiveIndicators}~: risque d'échec prévisionnel, mauvais choix de filière, risque de décrochage.
    \item \texttt{successForecast}~: taux de réussite académique prévisionnelle, trajectoire optimale.
  \end{itemize}
\end{enumerate} \\
\hline
\rowcolor{rowgray}Scénario alternatif &
\begin{itemize}
  \item L'utilisateur sélectionné n'est pas un étudiant $\rightarrow$ \texttt{abort(403)}.
\end{itemize} \\
\hline
Composants impliqués & \texttt{CounselorController::showStudent()}, \texttt{User}, \texttt{Profile}, \texttt{TestAttempt}, \texttt{Appointment} (Models), vue \texttt{counselor/student-profile.blade.php} \\
\hline
\end{tabular}

\vspace{1em}

%-------------------------------------------------------------
\subsection{Rédiger des observations de suivi}
%-------------------------------------------------------------
\ucheader{Rédiger des observations de suivi et mettre à jour le profil}{Conseiller d'orientation}

\noindent
\begin{tabular}{|L|R|}
\hline
\rowcolor{capblue}\textcolor{white}{\textbf{Champ}} & \textcolor{white}{\textbf{Description}} \\
\hline
\rowcolor{rowgray}Identifiant & UC-C03 \\
\hline
Cas d'utilisation & Rédiger des observations de suivi \\
\hline
\rowcolor{rowgray}Acteur & Conseiller d'orientation \\
\hline
Précondition & Dossier étudiant ouvert. \\
\hline
\rowcolor{rowgray}Postcondition & Observation (\texttt{counselor\_observations}), plan de coaching (\texttt{coaching\_plan}) et statut du dossier (\texttt{status}) mis à jour dans \texttt{Profile}. \\
\hline
Scénario nominal &
\begin{enumerate}
  \item Le conseiller remplit le formulaire de suivi (observations, plan de coaching, statut~: pending / ongoing / completed).
  \item Il soumet via \texttt{PATCH /counselor/students/\{student\}}.
  \item \texttt{CounselorController::updateProfile()} valide les données et met à jour le profil.
  \item La mise à jour apparaît dans la timeline CRM de l'étudiant.
\end{enumerate} \\
\hline
\rowcolor{rowgray}Scénario alternatif &
\begin{itemize}
  \item Statut invalide $\rightarrow$ La validation Laravel rejette la requête.
\end{itemize} \\
\hline
Composants impliqués & \texttt{CounselorController::updateProfile()}, \texttt{Profile} (Model) \\
\hline
\end{tabular}

\vspace{1em}

%-------------------------------------------------------------
\subsection{Planifier un rendez-vous d'orientation}
%-------------------------------------------------------------
\ucheader{Planifier un rendez-vous d'orientation}{Conseiller d'orientation}

\noindent
\begin{tabular}{|L|R|}
\hline
\rowcolor{capblue}\textcolor{white}{\textbf{Champ}} & \textcolor{white}{\textbf{Description}} \\
\hline
\rowcolor{rowgray}Identifiant & UC-C04 \\
\hline
Cas d'utilisation & Planifier un rendez-vous d'orientation \\
\hline
\rowcolor{rowgray}Acteur & Conseiller d'orientation \\
\hline
Précondition & Conseiller authentifié. \\
\hline
\rowcolor{rowgray}Postcondition & Rendez-vous enregistré dans \texttt{appointments} avec statut \texttt{scheduled}. \\
\hline
Scénario nominal &
\begin{enumerate}
  \item Le conseiller accède au planificateur depuis le dossier d'un étudiant.
  \item Il saisit la date/heure (\texttt{scheduled\_at}) et les notes du rendez-vous.
  \item Il soumet via \texttt{POST /counselor/students/\{student\}/appointments}.
  \item \texttt{CounselorController::storeAppointment()} valide que la date est dans le futur (\texttt{after:now}) et crée le rendez-vous.
\end{enumerate} \\
\hline
\rowcolor{rowgray}Scénario alternatif &
\begin{itemize}
  \item Date passée $\rightarrow$ La validation bloque la création avec un message d'erreur.
\end{itemize} \\
\hline
Composants impliqués & \texttt{CounselorController::storeAppointment()}, \texttt{Appointment} (Model) \\
\hline
\end{tabular}

\vspace{1em}

%-------------------------------------------------------------
\subsection{Valider ou ajuster le matching étudiant-filière}
%-------------------------------------------------------------
\ucheader{Valider / Ajuster le matching étudiant-filière (Homologation)}{Conseiller d'orientation}

\noindent
\begin{tabular}{|L|R|}
\hline
\rowcolor{capblue}\textcolor{white}{\textbf{Champ}} & \textcolor{white}{\textbf{Description}} \\
\hline
\rowcolor{rowgray}Identifiant & UC-C05 \\
\hline
Cas d'utilisation & Valider ou ajuster le matching étudiant-filière \\
\hline
\rowcolor{rowgray}Acteur & Conseiller d'orientation \\
\hline
Précondition & Recommandations SIAEPI générées et dossier étudiant ouvert. \\
\hline
\rowcolor{rowgray}Postcondition & Champ \texttt{manual\_match\_approved} mis à jour dans \texttt{Profile}~; \texttt{counselor\_observations} enrichi avec la justification horodatée. \\
\hline
Scénario nominal &
\begin{enumerate}
  \item Le conseiller sélectionne une action~: Approuver, Rejeter, Modifier la trajectoire ou Changer la filière.
  \item Il rédige une justification obligatoire (min. 10 caractères).
  \item Il soumet via \texttt{POST /counselor/students/\{student\}/approve-match}.
  \item \texttt{CounselorController::approveMatch()} applique la décision~:
  \begin{itemize}
    \item \textit{approve} $\rightarrow$ \texttt{manual\_match\_approved = true}, observation enregistrée.
    \item \textit{reject} $\rightarrow$ \texttt{manual\_match\_approved = false}, préfixe \texttt{[REJET IA]} ajouté.
    \item \textit{modify\_trajectory} / \textit{change\_field} $\rightarrow$ Mise à jour des intérêts et observations.
  \end{itemize}
  \item L'historique décisionnel (\texttt{decisionHistory}) est mis à jour.
\end{enumerate} \\
\hline
\rowcolor{rowgray}Scénario alternatif &
\begin{itemize}
  \item Justification vide ou trop courte ($<$ 10 caractères) $\rightarrow$ Validation Laravel bloque l'envoi.
\end{itemize} \\
\hline
Composants impliqués & \texttt{CounselorController::approveMatch()}, \texttt{Profile} (Model) \\
\hline
\end{tabular}

\vspace{1em}

%-------------------------------------------------------------
\subsection{Envoyer une communication omnicanale}
%-------------------------------------------------------------
\ucheader{Envoyer une communication omnicanale à un étudiant}{Conseiller d'orientation}

\noindent
\begin{tabular}{|L|R|}
\hline
\rowcolor{capblue}\textcolor{white}{\textbf{Champ}} & \textcolor{white}{\textbf{Description}} \\
\hline
\rowcolor{rowgray}Identifiant & UC-C06 \\
\hline
Cas d'utilisation & Envoyer une communication omnicanale \\
\hline
\rowcolor{rowgray}Acteur & Conseiller d'orientation \\
\hline
Précondition & Conseiller authentifié, dossier d'un étudiant ouvert. \\
\hline
\rowcolor{rowgray}Postcondition & Message envoyé via le canal sélectionné (Chat, Email, Notification Push, SMS) et enregistré dans le journal de communication CRM. \\
\hline
Scénario nominal &
\begin{enumerate}
  \item Le conseiller sélectionne un canal~: chat interne, email, notification push ou SMS.
  \item Il choisit un template (convocation, relance, conseil, alerte) ou rédige librement.
  \item Il soumet via \texttt{POST /counselor/students/\{student\}/send-message}.
  \item \texttt{CounselorController::sendMessage()} valide et enregistre le message dans la session CRM.
  \item Le journal de communication est mis à jour en temps réel.
\end{enumerate} \\
\hline
\rowcolor{rowgray}Scénario alternatif &
\begin{itemize}
  \item Corps du message vide ou trop court ($<$ 5 caractères) $\rightarrow$ Validation bloque l'envoi.
\end{itemize} \\
\hline
Composants impliqués & \texttt{CounselorController::sendMessage()}, \texttt{Message} (Model) \\
\hline
\end{tabular}

\vspace{1em}

%-------------------------------------------------------------
\subsection{Vérifier son statut d'approbation et consulter ses justificatifs}
%-------------------------------------------------------------
\ucheader{Vérifier son statut et consulter ses justificatifs}{Conseiller en attente / rejeté}

\noindent
\begin{tabular}{|L|R|}
\hline
\rowcolor{capblue}\textcolor{white}{\textbf{Champ}} & \textcolor{white}{\textbf{Description}} \\
\hline
\rowcolor{rowgray}Identifiant & UC-C07 \\
\hline
Cas d'utilisation & Vérifier son statut d'approbation et consulter ses justificatifs \\
\hline
\rowcolor{rowgray}Acteur & Conseiller d'orientation (en attente ou refusé) \\
\hline
Précondition & Conseiller authentifié ayant le statut \texttt{PENDING\_APPROVAL} ou \texttt{REJECTED}. \\
\hline
\rowcolor{rowgray}Postcondition & Redirection stricte et blocage d'accès aux fonctionnalités du tableau de bord. Visualisation des informations transmises et motifs de refus éventuels. \\
\hline
Scénario nominal &
\begin{enumerate}
  \item Le conseiller tente d'accéder à son tableau de bord d'orientation.
  \item Le middleware \texttt{counselor.approved} intercepte la requête.
  \item Le système identifie que le statut est \texttt{PENDING\_APPROVAL} et le redirige vers \texttt{/counselor/pending}.
  \item Le système présente une barre d'avancement (Création du compte $\rightarrow$ Dossier soumis $\rightarrow$ Attente d'évaluation).
  \item Les détails du dossier soumis (nom, spécialité, expérience, CV PDF) sont affichés pour vérification.
\end{enumerate} \\
\hline
\rowcolor{rowgray}Scénario alternatif &
\begin{itemize}
  \item Statut \texttt{REJECTED} $\rightarrow$ Le système affiche une bannière d'alerte rouge avec le motif du refus fourni par l'administrateur, et invite à contacter le support.
  \item Le conseiller clique sur « Se déconnecter » $\rightarrow$ Clôture sécurisée de la session.
\end{itemize} \\
\hline
Composants impliqués & \texttt{CounselorController::pending()}, middleware \texttt{counselor.approved}, vue \texttt{counselor/pending.blade.php} \\
\hline
\end{tabular}

\newpage

%=============================================================
\section{Acteur : Administrateur}
%=============================================================

L'administrateur est responsable de la configuration globale de la plateforme~: gestion des utilisateurs et de leurs rôles, administration du questionnaire RIASEC, import des filières, consultation du journal d'audit et paramétrage des critères de référence de matching.

%-------------------------------------------------------------
\subsection{Consulter le tableau de bord d'administration}
%-------------------------------------------------------------
\ucheader{Consulter le tableau de bord d'administration}{Administrateur}

\noindent
\begin{tabular}{|L|R|}
\hline
\rowcolor{capblue}\textcolor{white}{\textbf{Champ}} & \textcolor{white}{\textbf{Description}} \\
\hline
\rowcolor{rowgray}Identifiant & UC-A01 \\
\hline
Cas d'utilisation & Consulter le tableau de bord d'administration \\
\hline
\rowcolor{rowgray}Acteur & Administrateur \\
\hline
Précondition & Accès authentifié avec le rôle \texttt{admin}. \\
\hline
\rowcolor{rowgray}Postcondition & Indicateurs globaux affichés~: nombre d'inscrits par rôle, tests terminés, étudiants actifs. \\
\hline
Scénario nominal &
\begin{enumerate}
  \item L'administrateur accède à \texttt{GET /admin} (route protégée par le middleware \texttt{role:admin}).
  \item \texttt{Admin/DashboardController::index()} agrège les statistiques globales de la plateforme.
  \item Les KPIs principaux sont affichés~: utilisateurs inscrits, étudiants actifs, conseillers, tests passés.
\end{enumerate} \\
\hline
\rowcolor{rowgray}Scénario alternatif &
\begin{itemize}
  \item Accès sans rôle admin $\rightarrow$ Redirection automatique par le middleware \texttt{role:admin}.
\end{itemize} \\
\hline
Composants impliqués & \texttt{Admin/DashboardController}, vue \texttt{admin/dashboard.blade.php} \\
\hline
\end{tabular}

\vspace{1em}

%-------------------------------------------------------------
\subsection{Gérer les comptes utilisateurs et les rôles}
%-------------------------------------------------------------
\ucheader{Gérer les comptes utilisateurs (CRUD + Rôles)}{Administrateur}

\noindent
\begin{tabular}{|L|R|}
\hline
\rowcolor{capblue}\textcolor{white}{\textbf{Champ}} & \textcolor{white}{\textbf{Description}} \\
\hline
\rowcolor{rowgray}Identifiant & UC-A02 \\
\hline
Cas d'utilisation & Gérer les comptes utilisateurs et les rôles \\
\hline
\rowcolor{rowgray}Acteur & Administrateur \\
\hline
Précondition & Administrateur authentifié. \\
\hline
\rowcolor{rowgray}Postcondition & Statut, rôle ou validité d'un compte utilisateur mis à jour~; session active de l'utilisateur ciblé invalidée. \\
\hline
Scénario nominal &
\begin{enumerate}
  \item L'administrateur ouvre « Gestion des utilisateurs » via \texttt{GET /admin/users}.
  \item Le système affiche la table paginée de tous les comptes.
  \item L'administrateur peut~:
  \begin{itemize}
    \item Bloquer/Débloquer~: \texttt{POST /admin/users/block/\{user\}} (toggle \texttt{is\_blocked}).
    \item Promouvoir~: \texttt{POST /admin/users/promote/\{user\}} (étudiant $\rightarrow$ conseiller $\rightarrow$ admin).
    \item Rétrograder~: \texttt{POST /admin/users/demote/\{user\}}.
    \item Supprimer~: \texttt{DELETE /admin/users/\{user\}}.
  \end{itemize}
  \item Le système applique le changement et met à jour le champ \texttt{role} dans \texttt{users}.
\end{enumerate} \\
\hline
\rowcolor{rowgray}Scénario alternatif &
\begin{itemize}
  \item Tentative de blocage de son propre compte $\rightarrow$ Le système interdit l'action (protection dans \texttt{UserController}).
\end{itemize} \\
\hline
Composants impliqués & \texttt{Admin/UserController}, \texttt{User} (Model), vue \texttt{admin/users/index.blade.php} \\
\hline
\end{tabular}

\vspace{1em}

%-------------------------------------------------------------
\subsection{Gérer les questions du test RIASEC (CRUD)}
%-------------------------------------------------------------
\ucheader{Gérer les questions du test RIASEC (CRUD complet)}{Administrateur}

\noindent
\begin{tabular}{|L|R|}
\hline
\rowcolor{capblue}\textcolor{white}{\textbf{Champ}} & \textcolor{white}{\textbf{Description}} \\
\hline
\rowcolor{rowgray}Identifiant & UC-A03 \\
\hline
Cas d'utilisation & Gérer les questions du test RIASEC \\
\hline
\rowcolor{rowgray}Acteur & Administrateur \\
\hline
Précondition & Accès au module de test RIASEC. \\
\hline
\rowcolor{rowgray}Postcondition & Question créée, modifiée, désactivée ou supprimée dans le pool CAT (\texttt{riasec\_questions}). \\
\hline
Scénario nominal &
\begin{enumerate}
  \item L'administrateur ouvre \texttt{GET /admin/riasec/questions}.
  \item Le système affiche toutes les questions avec~: énoncé, paramètres 2PL (difficulté $b$, discrimination $a$), dimension Holland ciblée, catégorie GATB, statut inversé (\textit{reverse scored}), bloc, \textit{bacs cibles}.
  \item L'administrateur crée une question via \texttt{GET /admin/riasec/questions/create} $\rightarrow$ \texttt{POST /admin/riasec/questions}.
  \item Il peut éditer (\texttt{PUT}), supprimer (\texttt{DELETE}) ou basculer le statut (\texttt{toggle}).
  \item Il peut exporter toutes les données via \texttt{GET /admin/riasec/export} (CSV).
\end{enumerate} \\
\hline
\rowcolor{rowgray}Scénario alternatif &
\begin{itemize}
  \item Question déjà répondue par des étudiants $\rightarrow$ Recommandé de désactiver (toggle) plutôt que de supprimer, pour préserver l'intégrité des calculs historiques.
\end{itemize} \\
\hline
Composants impliqués & \texttt{Admin/RiasecAdminController}, \texttt{QuestionRiasec} (Model), vues \texttt{admin/riasec/} \\
\hline
\end{tabular}

\vspace{1em}

%-------------------------------------------------------------
\subsection{Importer des filières depuis Excel}
%-------------------------------------------------------------
\ucheader{Importer les filières et formations (Excel)}{Administrateur}

\noindent
\begin{tabular}{|L|R|}
\hline
\rowcolor{capblue}\textcolor{white}{\textbf{Champ}} & \textcolor{white}{\textbf{Description}} \\
\hline
\rowcolor{rowgray}Identifiant & UC-A04 \\
\hline
Cas d'utilisation & Importer des filières depuis un fichier Excel \\
\hline
\rowcolor{rowgray}Acteur & Administrateur \\
\hline
Précondition & Fichier Excel contenant les données de filières (SDO, RIASEC, sections BAC acceptées) prêt. \\
\hline
\rowcolor{rowgray}Postcondition & Table \texttt{filieres} mise à jour en masse par transaction SQL. \\
\hline
Scénario nominal &
\begin{enumerate}
  \item L'administrateur accède à \texttt{GET /admin/filieres/import}.
  \item Il sélectionne son fichier Excel et soumet via \texttt{POST /admin/filieres/import}.
  \item \texttt{FiliereImportController::store()} parse et valide le fichier (codes RIASEC, SDO, sections BAC).
  \item Les filières valides sont insérées en base par transaction SQL globale.
  \item Un rapport de résultat est affiché (lignes importées / lignes rejetées).
\end{enumerate} \\
\hline
\rowcolor{rowgray}Scénario alternatif &
\begin{itemize}
  \item Schéma invalide ou valeurs corrompues $\rightarrow$ Annulation complète (\textit{rollback}), rapport d'erreurs précisant les lignes défaillantes.
  \item Suppression par catégorie $\rightarrow$ \texttt{DELETE /admin/filieres/import/\{categorie\}}.
\end{itemize} \\
\hline
Composants impliqués & \texttt{Admin/FiliereImportController}, \texttt{Filiere}, \texttt{Specialite} (Models) \\
\hline
\end{tabular}

\vspace{1em}

%-------------------------------------------------------------
\subsection{Consulter le journal d'audit et de sécurité}
%-------------------------------------------------------------
\ucheader{Consulter le journal d'audit et de sécurité}{Administrateur}

\noindent
\begin{tabular}{|L|R|}
\hline
\rowcolor{capblue}\textcolor{white}{\textbf{Champ}} & \textcolor{white}{\textbf{Description}} \\
\hline
\rowcolor{rowgray}Identifiant & UC-A05 \\
\hline
Cas d'utilisation & Consulter le journal d'audit et de sécurité \\
\hline
\rowcolor{rowgray}Acteur & Administrateur \\
\hline
Précondition & Administrateur authentifié. \\
\hline
\rowcolor{rowgray}Postcondition & Journal des actions critiques affiché (promotions, connexions échouées, suppressions, modifications). \\
\hline
Scénario nominal &
\begin{enumerate}
  \item L'administrateur accède à \texttt{GET /admin/audit}.
  \item \texttt{AuditController::index()} récupère les entrées de \texttt{audit\_logs}~: action, utilisateur, IP, horodatage, détails.
  \item L'administrateur filtre les logs par type d'événement ou période.
\end{enumerate} \\
\hline
\rowcolor{rowgray}Scénario alternatif &
\begin{itemize}
  \item Journal volumineux $\rightarrow$ Pagination stricte et export CSV possible.
\end{itemize} \\
\hline
Composants impliqués & \texttt{Admin/AuditController}, \texttt{AuditLog} (Model), vue \texttt{admin/audit.blade.php} \\
\hline
\end{tabular}

\vspace{1em}

%-------------------------------------------------------------
\subsection{Gérer les critères et références de matching}
%-------------------------------------------------------------
\ucheader{Gérer les critères et références de matching (SIAEPI)}{Administrateur}

\noindent
\begin{tabular}{|L|R|}
\hline
\rowcolor{capblue}\textcolor{white}{\textbf{Champ}} & \textcolor{white}{\textbf{Description}} \\
\hline
\rowcolor{rowgray}Identifiant & UC-A06 \\
\hline
Cas d'utilisation & Gérer les critères et références de matching \\
\hline
\rowcolor{rowgray}Acteur & Administrateur \\
\hline
Précondition & Accès aux paramètres de configuration du moteur SIAEPI. \\
\hline
\rowcolor{rowgray}Postcondition & Sections de référence (\texttt{reference\_sections}) et critères (\texttt{reference\_criteria}) mis à jour pour les prochains calculs de matching. \\
\hline
Scénario nominal &
\begin{enumerate}
  \item L'administrateur ouvre \texttt{GET /admin/references}.
  \item Il consulte et modifie les sections et critères de référence.
  \item Il peut créer une section (\texttt{POST /admin/references}), un critère (\texttt{POST /admin/references/criteria}) ou les supprimer.
  \item Les nouveaux paramètres s'appliquent aux prochains calculs SIAEPI.
\end{enumerate} \\
\hline
\rowcolor{rowgray}Scénario alternatif &
\begin{itemize}
  \item Données incomplètes $\rightarrow$ La validation Laravel rejette la requête.
\end{itemize} \\
\hline
Composants impliqués & \texttt{Admin/ReferenceController}, \texttt{ReferenceSection}, \texttt{ReferenceCriterion} (Models), vue \texttt{admin/references.blade.php} \\
\hline
\end{tabular}

\vspace{1em}

%-------------------------------------------------------------
\subsection{Gérer les messages de contact}
%-------------------------------------------------------------
\ucheader{Gérer les messages de contact visiteurs}{Administrateur}

\noindent
\begin{tabular}{|L|R|}
\hline
\rowcolor{capblue}\textcolor{white}{\textbf{Champ}} & \textcolor{white}{\textbf{Description}} \\
\hline
\rowcolor{rowgray}Identifiant & UC-A07 \\
\hline
Cas d'utilisation & Gérer les messages de contact visiteurs \\
\hline
\rowcolor{rowgray}Acteur & Administrateur \\
\hline
Précondition & Administrateur authentifié. \\
\hline
\rowcolor{rowgray}Postcondition & Messages traités, lus ou supprimés. Compteur de notifications mis à jour. \\
\hline
Scénario nominal &
\begin{enumerate}
  \item L'administrateur accède à \texttt{GET /admin/contacts}.
  \item Le système liste les messages reçus de la landing page publique.
  \item L'administrateur consulte le détail d'un message via \texttt{GET /admin/contacts/\{id\}}.
  \item Il peut supprimer un message via \texttt{DELETE /admin/contacts/\{id\}}.
  \item Le compteur de notifications non lues est mis à jour via \texttt{GET /admin/notifications/count}.
\end{enumerate} \\
\hline
\rowcolor{rowgray}Scénario alternatif &
\begin{itemize}
  \item Aucun message $\rightarrow$ Affichage d'un état vide.
\end{itemize} \\
\hline
Composants impliqués & \texttt{ContactController}, \texttt{Contact} (Model), vue \texttt{admin/contacts/} \\
\hline
\end{tabular}

\vspace{1em}

%-------------------------------------------------------------
\subsection{Valider ou refuser les candidatures des conseillers}
%-------------------------------------------------------------
\ucheader{Valider ou refuser les candidatures des conseillers}{Administrateur}

\noindent
\begin{tabular}{|L|R|}
\hline
\rowcolor{capblue}\textcolor{white}{\textbf{Champ}} & \textcolor{white}{\textbf{Description}} \\
\hline
\rowcolor{rowgray}Identifiant & UC-A08 \\
\hline
Cas d'utilisation & Valider ou refuser les candidatures des conseillers \\
\hline
\rowcolor{rowgray}Acteur & Administrateur \\
\hline
Précondition & Administrateur authentifié. \\
\hline
\rowcolor{rowgray}Postcondition & Conseiller approuvé (rôle mis à jour à \texttt{counselor}, statut à \texttt{APPROVED}) ou rejeté (statut à \texttt{REJECTED} avec motif). \\
\hline
Scénario nominal &
\begin{enumerate}
  \item L'administrateur accède à « Candidatures Conseillers » via \texttt{GET /admin/counselors}.
  \item Le système affiche le tableau des demandes en attente (nom, spécialité, expérience, justificatif CV).
  \item L'administrateur clique sur « Approuver ».
  \item Le système met à jour le rôle de l'utilisateur à \texttt{counselor}, son statut à \texttt{APPROVED}, associe l'ID de l'admin et enregistre la date.
  \item L'utilisateur est notifié et peut désormais accéder aux fonctionnalités CRM du dashboard conseiller.
\end{enumerate} \\
\hline
\rowcolor{rowgray}Scénario alternatif &
\begin{itemize}
  \item L'administrateur clique sur « Refuser » $\rightarrow$ Une boîte modale s'ouvre pour saisir obligatoirement un motif de refus.
  \item Il saisit le motif de refus et valide $\rightarrow$ Le statut passe à \texttt{REJECTED} et le motif est persisté dans le profil.
\end{itemize} \\
\hline
Composants impliqués & \texttt{Admin/AdminCounselorController}, \texttt{User} (Model), \texttt{CounselorProfile} (Model), vue \texttt{admin/counselors/index.blade.php} \\
\hline
\end{tabular}

\vspace{2em}
\begin{center}
{\color{capblue}\rule{\linewidth}{1pt}}\\[0.5em]
{\small\textbf{CapAvenir} --- Cahier des Cas d'Utilisation --- 28 cas documentés sur 4 acteurs}\\[0.2em]
{\small Stack~: Laravel 10 $\cdot$ PHP 8.2 $\cdot$ Moteur SIAEPI v5.0 $\cdot$ CAT 2PL $\cdot$ API Gemini}\\[0.3em]
{\color{capblue}\rule{\linewidth}{1pt}}
\end{center}

\end{document}
